\subsection{Maschke's Theorem}

THEOREM 1. --- Suppose that the group G is finite. Let M be a K{[}G{]}-module, and let N be a K{[}G{]}-submodule of M. Suppose that N is a direct factor of the K-module M and that \textbar G\textbar{} is invertible in K. Then N is a direct factor of the K{[}G{]}-module M.

Let p be a K-linear projector in M with image N. We define an endomorphism q of the K-module M by setting

\[
q (x) = | \mathrm{G} | ^ {- 1} \sum_ {g \in \mathrm{G}} g p (g ^ {- 1} x)
\]

for every \(x \in M\) . Since N is stable under the action of G and p induces the identity on N, we see that q sends M into N and induces the identity on N.

The K-module M is therefore the direct sum of the image N of q and the kernel of q. We have \(g q(x) = q(gx)\) for every \(x \in M\) and \(g \in G\) , so that the kernel of q is a K{[}G{]}-submodule of M. This proves that N is a direct factor of the K{[}G{]}-module M.

COROLLARY 1 (Maschke). --- Suppose that the group G is finite and that K is a commutative field. The algebra K{[}G{]} is semisimple if and only if the element \textbar G\textbar{} of the field K is not zero.

Suppose \(|\mathrm{G}| \neq 0\) . By Theorem 1, every K{[}G{]}-submodule of \(\mathrm{K}[\mathrm{G}]_s\) is a direct factor. Hence K{[}G{]} is semisimple.

Conversely, suppose that \(|G|\) is zero, and denote by \(\varepsilon\) the element \(\sum_{g\in G} g\) of the center of K{[}G{]}. We have \(\varepsilon \neq 0\) but \(\varepsilon^{2} = |G|\varepsilon = 0\) , so K{[}G{]} is not semisimple (VIII, p.~153, Proposition 3, a) and VIII, p.~157, Remark 1).

COROLLARY 2. --- Suppose that the group G is finite and that \(|G|\) is invertible in K. An exact sequence of K{[}G{]}-modules splits if and only if it splits as an exact sequence of K-modules.

Given an exact sequence of K{[}G{]}-modules

\[
0 \longrightarrow \mathrm{M} ^ {\prime} \stackrel {f} {\longrightarrow} \mathrm{M} \longrightarrow \mathrm{M} ^ {\prime \prime} \longrightarrow 0,
\]

it suffices to apply Theorem 1 to the image of the morphism f.

COROLLARY 3. --- Suppose that the group G is finite and that \(|G|\) is invertible in K. A K{[}G{]}-module is projective if and only if it is projective as a K-module.

Let P be a K{[}G{]}-module. If P is a direct factor of a free K{[}G{]}-module M, then it is a fortiori a direct factor of the free K-module M. The converse follows from Corollary 2, in view of II, §2, No.~2, p.~231, Proposition 4, d).

COROLLARY 4. --- a) Suppose that the group G is finite and that K is a commutative field of characteristic 0. Two finite-dimensional linear representations of G are isomorphic if and only if they have the same character.

\begin{enumerate}
\def\labelenumi{\alph{enumi})}
\setcounter{enumi}{1}
\tightlist
\item
  Suppose that K is a perfect field of characteristic a prime number p and that the group G is finite, of cardinal prime to p.~The character of a finite-dimensional linear representation of G is zero if and only if this representation is isomorphic to the direct sum of p pairwise isomorphic representations.
\end{enumerate}

Under the assumptions of the corollary, every finite-dimensional linear representation of G is semisimple. The corollary then follows from the corollary of VIII, p.~384.

COROLLARY 5. --- Suppose that the group G is finite and that K is a commutative field in which \(|G| \neq 0\) . Let \(\pi\) and \(\pi'\) be finite-dimensional linear representations of G. Then \(\pi\) and \(\pi'\) are isomorphic if and only if for every g in G, the endomorphisms \(\pi(g)\) and \(\pi'(g)\) have the same characteristic polynomial.

This follows from Corollary 1 of VIII, p.~386.

